\section{Archimedean decay I: logarithmic energy and positive arguments}\label{sec:energy}

This section and the next bound $\log|\Qt_n(\xi)|$ from above in terms of a comparison measure; with the measures of \cref{app:certificates} this gives \cref{thm:arch} (\cref{sec:proof-arch}). The two ingredients common to both are a regularized energy inequality (\cref{sec:energy-inequality}) and Stirling's formula for the normalization (\cref{lem:stirling}). For positive $\lambda$ we then use Heine's formula for the derivative functional of \cref{lem:series} directly (\cref{prop:positive-bound}). The statements are candidates in the sense of \cref{sec:intro-status}. The structure of the argument is that of \cite[Section~6]{Litsey2026Li2}: Heine's formula for a functional with derivatives, in which a site carries at most two variables (Lemma~6.2 there), a comparison with the energy of a measure, and a bound with a level $\theta$ and a penalty integral (Theorem~6.1 there, for the lattice). The energy inequality is proved here by a different, elementary route (\cref{lem:neutral-energy,lem:energy-comparison}).

\subsection{An energy inequality}\label{sec:energy-inequality}
For $\eps>0$ and $z,w\in\C$ put $\mathcal K_\eps(z,w)=\tfrac12\log(|z-w|^2+\eps^2)$.

\begin{lemma}[neutral charges]\label{lem:neutral-energy}
For $\eps>0$, points $z_1,\dots,z_m\in\C$ and real numbers $b_1,\dots,b_m$ with $\sum_ib_i=0$,
\[
\sum_{i,j}b_ib_j\,\mathcal K_\eps(z_i,z_j)\le0 .
\]
\end{lemma}

\begin{proof}
For $\theta>0$ the matrix $[e^{-\theta|z_i-z_j|^2}]$ is positive semidefinite: writing $z=x+\mathrm iy$, it equals the congruence by the diagonal matrix $\mathrm{diag}(e^{-\theta|z_i|^2})$ of the matrix $[e^{2\theta(x_ix_j+y_iy_j)}]$, which is a convergent sum of the Hadamard powers $[(x_ix_j+y_iy_j)^k]$ with nonnegative coefficients, each positive semidefinite by Schur's product theorem. Since $(s+\varrho^2)^{-1}=\int_0^\infty e^{-\theta s}e^{-\theta\varrho^2}\dd\theta$ for $s>0$, the matrix $[(s+|z_i-z_j|^2)^{-1}]$ is positive semidefinite. Hence $\varphi(s)=\sum_{i,j}b_ib_j\log(s+|z_i-z_j|^2)$ is nondecreasing on $(0,\infty)$. As $\sum b_i=0$, $\varphi(s)=\sum_{i,j}b_ib_j\log(1+|z_i-z_j|^2/s)\to0$ when $s\to\infty$. So $\varphi(\eps^2)\le0$.
\end{proof}

A \emph{comparison measure} is a positive Borel measure $\sigma$ on $\C$ of total mass $1$, supported on a compact subset of a finite union of $m_\sigma$ straight lines, and absolutely continuous with respect to arc length on each line with density at most $\varrho_\sigma$. Its potential and energy are
\[
U_\sigma(z)=\int\log|z-w|\dd\sigma(w),\qquad E(\sigma)=\int U_\sigma\dd\sigma .
\]
The potential is continuous on $\C$, and $E(\sigma)$ is finite.

\begin{lemma}[energy comparison]\label{lem:energy-comparison}
Let $\sigma$ be a comparison measure. For every $h\ge2$ and all $z_1,\dots,z_h\in\C$,
\[
\sum_{\substack{i\ne j\\ z_i\ne z_j}}\log|z_i-z_j|\ \le\ -h^2E(\sigma)+2h\sum_{i=1}^hU_\sigma(z_i)+2\,n_c\log h+2\pi m_\sigma\varrho_\sigma,
\]
where $n_c$ is the number of ordered pairs $(i,j)$, including $i=j$, with $z_i=z_j$.
\end{lemma}

\begin{proof}
Let $\eps=h^{-2}$, $U^\eps(z)=\int \mathcal K_\eps(z,w)\dd\sigma(w)$ and $E^\eps=\int U^\eps\dd\sigma$. Approximate $\sigma$ weakly by finite positive point measures $\sigma_k$ of mass one supported in $\operatorname{supp}\sigma$ (Riemann sums on each line). Since $\mathcal K_\eps$ is continuous, \cref{lem:neutral-energy} applied to $\mu-\sigma_k$, with $\mu=h^{-1}\sum_i\delta_{z_i}$, gives in the limit $\iint \mathcal K_\eps\dd(\mu-\sigma)\dd(\mu-\sigma)\le0$, that is
\[
\sum_{i,j}\mathcal K_\eps(z_i,z_j)\le2h\sum_iU^\eps(z_i)-h^2E^\eps .
\]
On the left, $\mathcal K_\eps(z,w)\ge\log|z-w|$ for $z\ne w$ and $\mathcal K_\eps(z,z)=\log\eps=-2\log h$. On the right, $E^\eps\ge E(\sigma)$, and $0\le U^\eps(z)-U_\sigma(z)\le\pi m_\sigma\varrho_\sigma\eps$: on each line, parametrized by arc length $u$, with $u_0$ the parameter of the orthogonal projection of $z$, we have $|z-w|\ge|u-u_0|$ and $\int_\R\frac12\log\bigl(1+\eps^2(u-u_0)^{-2}\bigr)\dd u=\pi\eps$. Combining these estimates proves the lemma.
\end{proof}

\begin{lemma}[normalization]\label{lem:stirling}
As $n\to\infty$,
\[
\log\frac{S_n^h}{F_n}=(d-1)h^2\log h+\Cst h^2+O(h\log h),\qquad \Cst=\kappa\log\kappa-3a\log a+\tfrac32-d,
\]
with $0\log0=0$.
\end{lemma}

\begin{proof}
By Stirling's formula $\log m!=m\log m-m+O(\log(m+1))$, and $\sum_{i<h}(i\log i-i)=\frac{h^2}2\log h-\frac34h^2+O(h\log h)$. Hence $h\log K!-3h\log A!=(\kappa-3a)h^2\log h+(\kappa\log\kappa-\kappa-3a\log a+3a)h^2+O(h\log h)$ and $\log F_n=h^2\log h-\frac32h^2+O(h\log h)$.
\end{proof}

The constants are $\Cst=2.3979184$, $0.5075059$, $2.5190484$ and $2.2176686$ for $\lambda=2/11$, $2/9$, $-2/9$ and $-2/7$.

\subsection{Heine's formula with derivatives}\label{sec:heine}
For a linear functional $\ell$ on $\C[x]$ let $\ell^{\otimes h}$ be the linear functional on $\C[x_1,\dots,x_h]$ with $\ell^{\otimes h}(\prod_ix_i^{k_i})=\prod_i\ell(x^{k_i})$, and let $\Delta(x_1,\dots,x_h)=\prod_{i<j}(x_j-x_i)$.

\begin{lemma}[Heine]\label{lem:heine}
$\det_{0\le i,j<h}\ell(x^{i+j})=\frac1{h!}\,\ell^{\otimes h}(\Delta^2)$.
\end{lemma}

\begin{proof}
Expanding $\Delta=\det[x_i^{j}]$ over permutations,
\[
\Delta=\sum_\pi\operatorname{sgn}(\pi)\prod_ix_i^{\pi(i)},\qquad \ell^{\otimes h}(\Delta^2)=\sum_{\pi,\pi'}\operatorname{sgn}(\pi\pi')\prod_i\ell\bigl(x^{\pi(i)+\pi'(i)}\bigr),
\]
and for each fixed $\pi$ the sum over $\pi'$ equals $\det[\ell(x^{i+j})]$ after reordering the rows.
\end{proof}

By \cref{lem:series}, the functional $\ell(\phi)=\Lambda_\xi(W_n\phi)$ on polynomials is
\[
\ell(\phi)=\sum_{x\ge1}\bigl(\alpha_x\phi(x)+\beta_x\phi'(x)\bigr),\qquad \alpha_x=\lambda^x\bigl(xW_n(x)\bigr)',\quad \beta_x=\lambda^xxW_n(x),
\]
since $(xW_n\phi)'=(xW_n)'\phi+xW_n\phi'$. The series converges absolutely for every polynomial $\phi$, so that $\ell^{\otimes h}=\sum_{x_1,\dots,x_h\ge1}\ell_{x_1}\otimes\dots\otimes\ell_{x_h}$ on polynomials, with $\ell_x(\phi)=\alpha_x\phi(x)+\beta_x\phi'(x)$, and
\begin{equation}\label{eq:heine-sum}
Q_n(\xi)=\frac1{h!}\sum_{\mathbf x\in\Z_{\ge1}^h}\mathfrak T(\mathbf x),\qquad \mathfrak T(\mathbf x)=(\ell_{x_1}\otimes\dots\otimes\ell_{x_h})(\Delta^2).
\end{equation}

\begin{lemma}[multiplicity]\label{lem:multiplicity}
There is a constant $C_0$, depending only on $\lambda$ and $(H,B,C)$, with the following property. If some value occurs at least three times among $x_1,\dots,x_h$, then $\mathfrak T(\mathbf x)=0$. Otherwise
\[
|\mathfrak T(\mathbf x)|\ \le\ (C_0h)^h\prod_{i=1}^h|\beta_{x_i}|\prod_{\substack{i\ne j\\ x_i\ne x_j}}|x_i-x_j| .
\]
\end{lemma}

\begin{proof}
For a site $s$ let $I_s=\{i:x_i=s\}$ and $m_s=|I_s|$. Write $\Delta^2=\Phi_{\mathrm{int}}\Phi_{\mathrm{ext}}$, where $\Phi_{\mathrm{int}}=\prod_s\Phi_s$ with $\Phi_s=\prod_{i<j\in I_s}(x_i-x_j)^2$, and $\Phi_{\mathrm{ext}}$ is the product of the factors $(x_i-x_j)^2$ with $x_i\ne x_j$. The functional $\bigotimes_i\ell_{x_i}$ is the sum over $D\subset\{1,\dots,h\}$ of $\prod_{i\notin D}\alpha_{x_i}\prod_{i\in D}\beta_{x_i}$ times evaluation at $\mathbf x$ of $\partial_D=\prod_{i\in D}\partial/\partial x_i$. The polynomial $\Phi_s$ is homogeneous of degree $m_s(m_s-1)$ in the differences $x_i-x_j$ ($i,j\in I_s$); any derivative of it of order at most $m_s$ in these variables vanishes at the diagonal point unless its degree $m_s(m_s-1)-(\text{order})$ is zero. Hence $\mathfrak T(\mathbf x)=0$ if some $m_s\ge3$, and if $m_s=2$ only the terms with $I_s\subset D$ in which both derivatives fall on $\Phi_s$ survive, with $\partial_{x_i}\partial_{x_j}(x_i-x_j)^2=-2$. Writing $\mathsf S_1$ for the indices at singly occupied sites,
\[
\mathfrak T(\mathbf x)=\prod_{m_s=2}\bigl(-2\beta_s^2\bigr)\sum_{D\subset\mathsf S_1}\prod_{i\in\mathsf S_1\setminus D}\alpha_{x_i}\prod_{i\in D}\beta_{x_i}\,(\partial_D\Phi_{\mathrm{ext}})(\mathbf x).
\]
By Leibniz's rule each derivative $\partial/\partial x_i$ falls on one of the factors of $\Phi_{\mathrm{ext}}$ that contain $x_i$, of which there are fewer than $h$; so $\partial_D\Phi_{\mathrm{ext}}$ is a sum of at most $h^{|D|}$ terms, in each of which every factor $(x_i-x_j)^2$ is differentiated at most twice (once in $x_i$ and once in $x_j$). The derivatives of $(x_i-x_j)^2$ are $\pm2(x_i-x_j)$ and $-2$. Since $\Phi_{\mathrm{ext}}$ contains only factors with $x_i\ne x_j$, and the $x_i$ are integers, $|x_i-x_j|\ge1$ for each of them, so each differentiated factor is at most $2(x_i-x_j)^2$ in absolute value. So $|\partial_D\Phi_{\mathrm{ext}}(\mathbf x)|\le(2h)^{|D|}|\Phi_{\mathrm{ext}}(\mathbf x)|$. Finally $\alpha_x/\beta_x=(xW_n)'/(xW_n)=x^{-1}+3\sum_{a\le A}(x+a)^{-1}-\sum_{l\le K}(x+l)^{-1}$, so $|\alpha_x|\le(4+4\log K)|\beta_x|$. Summing over $D$ gives the bound with $C_0h\ge2(4+4\log K+2h)$.
\end{proof}

\begin{lemma}[the weight on the integers]\label{lem:weight-positive}
For $y\ge0$ let
\[
V(y)=3(y+a)\log(y+a)-2y\log y-(y+\kappa)\log(y+\kappa)+d-ty ,
\]
with $0\log0=0$ and $t=-\log|\lambda|$. There is $C_1$ such that for all $n$ and all integers $x\ge1$, with $y=x/h$,
\[
\log|\beta_x|\le-dh\log h+hV(y)+C_1\log(h+x).
\]
\end{lemma}

\begin{proof}
We have $\log|\beta_x|=-tx+\log x+3\log D_A(x)-\log D_K(x)$. Since $s\mapsto\log(x+s)$ is increasing, $\log D_m(x)\le\int_1^{m+1}\log(x+s)\dd s$ and $\log D_m(x)\ge\int_0^m\log(x+s)\dd s$. Therefore
\[
3\log D_A(x)-\log D_K(x)\le3\int_0^A\log(x+s)\dd s-\int_0^K\log(x+s)\dd s+3\log(x+A+1).
\]
Substituting $x=hy$ and $s=hs'$, $\int_0^{\mu h}\log(x+s)\dd s=\mu h\log h+h\bigl(F(y+\mu)-F(y)\bigr)$ with $F(z)=z\log z-z$. Collecting terms, $3\int_0^A-\int_0^K=-dh\log h+h\bigl(V(y)+ty\bigr)$.
\end{proof}

\subsection{The bound for positive arguments}\label{sec:positive-bound}

\begin{proposition}[comparison bound, positive arguments]\label{prop:positive-bound}
Let $\lambda\in\{2/11,2/9\}$, let $\sigma$ be a comparison measure supported in $[0,\infty)$, write $U=U_\sigma$, $E=E(\sigma)$ and $\mathsf W=V+2U$. Then for every $\theta\in\R$,
\[
\limsup_{n\to\infty}\frac{\log|\Qt_n(\xi)|}{h^2}\ \le\ \Cst-E+\theta+2\int_0^\infty\bigl(\mathsf W(y)-\theta\bigr)_+\dd y .
\]
\end{proposition}

\begin{proof}
By \eqref{eq:heine-sum} and \cref{lem:multiplicity}, $|Q_n(\xi)|\le\sum_{\mathbf x}|\mathfrak T(\mathbf x)|$, the sum running over tuples in which no value occurs more than twice. Fix such a tuple, put $y_i=x_i/h$, and note that $n_c\le2h$. Since $\log h>0$ and there are at most $h^2$ ordered pairs,
\[
\sum_{\substack{i\ne j\\x_i\ne x_j}}\log|x_i-x_j|\le h^2\log h+\sum_{\substack{i\ne j\\y_i\ne y_j}}\log|y_i-y_j|\le h^2\log h-h^2E+2h\sum_iU(y_i)+O(h\log h)
\]
by \cref{lem:energy-comparison}. With \cref{lem:multiplicity,lem:weight-positive},
\begin{equation}\label{eq:term-bound}
\log|\mathfrak T(\mathbf x)|\le(1-d)h^2\log h-h^2E+h\sum_i\mathsf W(y_i)+C_1\sum_i\log(h+x_i)+O(h\log h).
\end{equation}
As $y\to\infty$, $V(y)=-ty-d\log y+O(1/y)$ and $U(y)=\log y+O(1/y)$, so there are $Y$ and $c_1>0$ with $\mathsf W(y)-\theta\le-c_1(y-Y)$ for $y\ge Y$; hence $\mathsf W-\theta\le(\mathsf W-\theta)_+-c_1(y-Y)_+$ everywhere. Since each value occurs at most twice among the $x_i$,
\[
\sum_i\mathsf W(y_i)\le h\theta+2\sum_{x\ge1}\bigl(\mathsf W(x/h)-\theta\bigr)_+-\frac{c_1}h\sum_i(x_i-Yh)_+ .
\]
Insert this into \eqref{eq:term-bound}. The resulting bound for $|\mathfrak T(\mathbf x)|$ was derived for the tuples in which no value occurs more than twice; it is a positive quantity defined for every $\mathbf x\in\Z_{\ge1}^h$, so we may sum it over the admissible tuples and then enlarge the sum to all of $\Z_{\ge1}^h$, where it factorizes. We obtain
\begin{multline*}
\log|Q_n(\xi)|\le(1-d)h^2\log h+h^2(\theta-E)+2h\sum_{x\ge1}\bigl(\mathsf W(x/h)-\theta\bigr)_+\\
+h\log\Bigl(\sum_{x\ge1}(h+x)^{C_1}e^{-c_1(x-Yh)_+}\Bigr)+O(h\log h).
\end{multline*}
The last sum is $O(h^{C_1+1})$. The function $(\mathsf W-\theta)_+$ is continuous on $[0,\infty)$ with compact support, so $h^{-1}\sum_{x\ge1}(\mathsf W(x/h)-\theta)_+\to\int_0^\infty(\mathsf W-\theta)_+$. Adding \cref{lem:stirling} proves the proposition.
\end{proof}

\begin{remark}[scale for $\lambda=2/11$]\label{rem:211-scale}
For $\lambda=2/11$ the certificate is formulated in the variable $s=x/n=4y$ and divides by $n^2=h^2/16$. The substitution $y=s/4$ transforms \cref{prop:positive-bound} into
\[
\limsup_{n\to\infty}\frac{\log|\Qt_n(\xi)|}{n^2}\le24\log6-16\log4+12-E'+4\theta'+2\int_0^\infty\bigl(\mathsf W'(s)-\theta'\bigr)_+\dd s,
\]
where $\sigma'$ is the image of $4\sigma$ under $y\mapsto4y$ (a measure of mass $4$), $E'=\iint\log|s-s'|\dd\sigma'\dd\sigma'$, $\mathsf W'=V'+2U_{\sigma'}$ and $V'(s)=3(s+1)\log(s+1)-2s\log s-(s+6)\log(s+6)+3-ts$. Indeed $E'=16E+16\log4$, $\mathsf W'(4y)=4\mathsf W(y)+5\log4$, and with $\theta'=4\theta+5\log4$ the right side equals $16$ times that of \cref{prop:positive-bound}.
\end{remark}
